\section{A finite-dimensional fixed-point theorem}
\label{sec:fixedpoint}

We prove the finite-dimensional fixed-point theorem of
Brouwer~\cite{Brouwer} used in the extraction step, including its
combinatorial ingredient.

\begin{theorem}[Finite-dimensional fixed point]\label{thm:fixedpoint}
Let \(K\) be a nonempty compact convex subset of a finite-dimensional
real vector space. Every continuous map \(F:K\to K\) has a fixed point.
\end{theorem}

\subsection{Triangulations with arbitrarily small mesh}

A \(d\)-simplex is the convex hull of \(d+1\) affinely independent
points; its faces are the convex hulls of subsets of its vertices.
A finite triangulation of a \(d\)-simplex consists of finitely many
\(d\)-simplices and all their faces, covering the original simplex,
such that two members intersect in a common face or not at all.
Its \emph{mesh} is the largest diameter of its members.
We first construct such triangulations with mesh tending to zero.

Let \(S\) have vertices \(a_0,\ldots,a_d\). For every nonempty
vertex set \(A\), write \(b_A=|A|^{-1}\sum_{a\in A}a\).
The barycentric subdivision of \(S\) has one simplex
\[
 \operatorname{conv}\{b_{A_0},\ldots,b_{A_d}\},\qquad
 A_k=\{a_{\pi(0)},\ldots,a_{\pi(k)}\},
\]
for each permutation \(\pi\) of \(\{0,\ldots,d\}\), together with
all faces of these simplices. The displayed barycenters are affinely
independent: their coordinate columns in the ordered vertices of
\(S\) form an invertible triangular matrix.

Here is a direct verification of the covering and intersection
properties. Write \(x=\sum_i\lambda_i a_i\), where
\(\lambda_i\ge0\) and \(\sum_i\lambda_i=1\), and order the
coordinates so that
\(\lambda_{\pi(0)}\ge\cdots\ge\lambda_{\pi(d)}\).
With \(\lambda_{\pi(d+1)}:=0\), put
\[
 t_k=(k+1)(\lambda_{\pi(k)}-\lambda_{\pi(k+1)}).
\]
Then
\[
 t_k\ge0,\quad \sum_{k=0}^d t_k=1,\quad
 x=\sum_{k=0}^d t_k b_{A_k}.
\]
Conversely, a convex combination of these barycenters has coordinates
in that order. Whenever \(t_k>0\), the prefix \(A_k\) is the set of
vertices whose coordinates are at least the corresponding positive
coordinate level. Thus the barycenters with positive coefficients
are determined by \(x\), independently of how tied coordinates are
ordered. Two subdivision simplices consequently intersect exactly
in the convex hull of their common vertices. Their faces inherit the
same property. On any face of \(S\), coordinates outside that face
are zero, so this construction restricts exactly to the barycentric
subdivision of that face. Subdividing each member of a triangulation
therefore gives another finite triangulation: the constructions agree
on shared faces, and simplices in different original members can
intersect only there.

For nested nonempty vertex sets \(A\subsetneq B\),
\[
 b_B=\frac{|A|}{|B|}b_A+
       \left(1-\frac{|A|}{|B|}\right)b_{B\setminus A},
 \qquad
 \|b_B-b_A\|\le\frac{d}{d+1}\operatorname{diam}(S).
\]
The diameter of a convex hull is at most the largest distance between
its vertices, by the triangle inequality applied to convex
combinations. Every new simplex thus has diameter at most
\(d/(d+1)\) times that of its parent. For \(d\ge1\), iterating
gives mesh tending to zero. The case \(d=0\) is immediate.

\subsection{Sperner's counting argument}

Write
\[
 \Delta_d=\{x\in\R^{d+1}:x_i\ge0\ (0\le i\le d),\quad
                                  \sum_{i=0}^d x_i=1\}.
\]
A triangulation restricts to a triangulation of each face of
\(\Delta_d\): intersecting a small simplex with a coordinate
hyperplane \(x_i=0\) selects precisely its vertices on that
hyperplane, since all coordinates are nonnegative. The intersections
of full dimension cover the face: their union is closed, while the
remaining finitely many lower-dimensional simplices cannot cover
any relatively open subset of the face. Every smaller member is a
face of a full-dimensional one, by the common-face property applied
at a relative interior point of the smaller member.

\begin{lemma}[Sperner~\cite{Sperner1928}]\label{lemma:sperner}
Label every vertex \(v\) of a finite triangulation of \(\Delta_d\)
by an integer \(\ell(v)\in\{0,\ldots,d\}\), with
\(v_{\ell(v)}>0\). The number of \(d\)-simplices whose vertex
labels are \(0,\ldots,d\), each occurring once, is odd.
\end{lemma}
\begin{proof}
For \(d\ge1\), we use the following elementary incidence fact.
A small \((d-1)\)-face belongs to two small \(d\)-simplices if
it is not contained in the boundary, and to one if it is.
To see this, choose a point in its relative interior. Any simplex containing that point
contains the entire face, by the common-face intersection property.
Finiteness allows a sufficiently small ball in the affine hull of
\(\Delta_d\), centered at the point, to avoid all other simplices.
Shrinking the ball further, each incident \(d\)-simplex occupies
exactly one half of this ball, on one side of the face's affine
hyperplane: the point lies on none of its other facets. Disjoint
interiors permit at most one simplex on either side, and the covering
property supplies all sides belonging to \(\Delta_d\). A boundary
\((d-1)\)-face lies in exactly one original facet, so precisely one
side belongs to \(\Delta_d\). Indeed, a coordinate vanishing at a
relative interior point vanishes on the whole face, and two distinct
original facets intersect in dimension at most \(d-2\).

Induct on \(d\). For \(d=0\), there is one vertex, labeled zero.
For \(d\ge1\), call a small \((d-1)\)-face \emph{special} if
its labels are exactly \(0,\ldots,d-1\). Count incidences between
small \(d\)-simplices and their special faces. A simplex with all
\(d+1\) labels contributes one incidence. A simplex containing
\(0,\ldots,d-1\) but not \(d\) has exactly one repeated label
and contributes two, by deleting either vertex with that label.
Every other simplex contributes zero. Modulo two, the incidence
count is therefore the number asserted to be odd.

By the incidence fact, the same count modulo two is the number of
special faces on the boundary. Such a face cannot lie in
\(x_i=0\) for \(i<d\), since then none of its vertices could
receive label \(i\). Hence these are exactly the fully labeled
top-dimensional simplices in the induced triangulation of
\(x_d=0\). The labels there satisfy the same rule in dimension
\(d-1\), and induction makes their number odd.
\end{proof}

\subsection{From a simplex to a compact convex set}

\begin{proof}[Proof of Theorem~\ref{thm:fixedpoint}]
First let \(F:\Delta_d\to\Delta_d\) be continuous. The case
\(d=0\) is immediate. In each of the successive barycentric
subdivisions, a vertex fixed by \(F\) already proves the claim.
Otherwise label each vertex \(v\) by an index \(i\) for which
\(v_i>F_i(v)\). Such an index exists because the two coordinate
sums are equal and \(v\ne F(v)\); it satisfies \(v_i>0\).
Lemma~\ref{lemma:sperner} supplies a fully labeled simplex in each
subdivision. Denote its vertex labeled \(i\) by \(v^{(m,i)}\),
where \(m\) is the subdivision index. By compactness, a subsequence
of \(v^{(m,0)}\) converges to some \(x\in\Delta_d\).
The mesh tends to zero, so all \(v^{(m,i)}\) converge to the same
\(x\) along that subsequence. Continuity gives
\(x_i\ge F_i(x)\) for every \(i\). Equality of the coordinate
sums forces equality in every coordinate, proving \(F(x)=x\).
Affine coordinates give the same result for every simplex.

Now equip the affine hull of \(K\) with Euclidean coordinates.
For any point \(x\) in this hull, compactness gives a nearest
point \(P(x)\in K\). It is unique: if distinct \(p,q\in K\)
both minimized the squared distance, convexity and
\[
 \left\|x-\frac{p+q}{2}\right\|^2
 =\frac{\|x-p\|^2+\|x-q\|^2}{2}
  -\frac{\|p-q\|^2}{4}
\]
would contradict minimality. If \(p=P(x)\) and \(z\in K\),
comparison with \(p+t(z-p)\), followed by division by \(t>0\)
and the limit \(t\downarrow0\), yields
\(\langle x-p,z-p\rangle\le0\).
Apply this with \(p=P(x),z=P(y)=q\), and with \(q=P(y),z=p\).
Adding the inequalities gives
\[
 \|p-q\|^2\le\langle x-y,p-q\rangle
              \le\|x-y\|\,\|p-q\|.
\]
Thus \(\|P(x)-P(y)\|\le\|x-y\|\), so \(P\) is continuous,
and \(P(x)=x\) for \(x\in K\).

If the affine hull has dimension zero, \(K\) is a singleton.
Otherwise identify it with \(\R^d\), \(d\ge1\), and choose
\(M>0\) so that \(K\subseteq[-M,M]^d\). The simplex
\[
 S=\{x\in\R^d:x_j\ge-M\ (1\le j\le d),\quad
                         \sum_{j=1}^d x_j\le dM\}
\]
contains that cube: its vertices are
\(a=(-M,\ldots,-M)\) and \(a+2dM e_j\), \(1\le j\le d\).
The continuous map \(F\circ P:S\to K\subseteq S\) has a fixed
point \(x\) by the simplex case. Its image lies in \(K\), so
\(x\in K\), \(P(x)=x\), and \(F(x)=x\), as required.
\end{proof}
